\section{User Management}


\subsection{User Registration Flows}

Lux.id supports multiple registration patterns:

\begin{enumerate}
    \item \textbf{Traditional}: Email/password with verification
    \item \textbf{Social}: OAuth providers (Google, GitHub, Twitter)
    \item \textbf{Wallet}: Direct blockchain wallet connection
    \item \textbf{Federated}: SAML/OIDC from enterprise IdP
    \item \textbf{Hybrid}: Wallet + email for recovery
\end{enumerate}

\subsection{Profile Management}

User profiles combine Web2 and Web3 attributes:

\begin{lstlisting}[language=json]
{
  "id": "user-123",
  "email": "user@lux.network",
  "wallets": [
    {
      "chain": "ethereum",
      "address": "0x742d35Cc...",
      "ens": "user.eth"
    },
    {
      "chain": "lux",
      "address": "lux1qpz3a..."
    }
  ],
  "credentials": {
    "webauthn": ["credentialId1", "credentialId2"],
    "totp": true
  },
  "attributes": {
    "kyc_verified": true,
    "validator_status": "active"
  }
}
\end{lstlisting}

\subsection{Role-Based Access Control (RBAC)}

RBAC implementation follows NIST model:

\begin{equation}
    \text{Permission} = \text{Role} \times \text{Operation} \times \text{Resource}
\end{equation}

Role hierarchy example:

\begin{lstlisting}
Admin
+-- Manager
|   +-- Validator
|   +-- Developer
+-- Auditor
    +-- ReadOnly
\end{lstlisting}

\subsection{Attribute-Based Access Control (ABAC)}

ABAC enables fine-grained access control:

\begin{lstlisting}[language=json]
{
  "policy": "allow",
  "conditions": [
    {"attribute": "department", "operator": "equals", "value": "trading"},
    {"attribute": "clearance", "operator": ">=", "value": 3},
    {"attribute": "nft.balance", "operator": ">", "value": 0}
  ],
  "resource": "/api/trading/advanced",
  "actions": ["read", "write"]
}
\end{lstlisting}

